% -----------------------------------------------------------------
\subsection{Functional Test Suite}
\label{subsec:crummy-test}

\code{CrummyTest} is a self-contained functional test suite invoked
from the application's right-click context menu
(\code{IDM\_BASICTEST}, labeled ``Run Basic Test'').  The suite
exercises the direct rendering path
(\cref{subsec:two-render-paths}): each test calls \code{Drain()}
first to ensure no pool workers remain active, then invokes
\code{CalcFractal(true)}, which writes results to \code{m\_CurIters}.

The top-level \code{TestAll()} method runs the full suite in sequence:
\begin{itemize}
  \item \code{TestWindowResize}---resizes the application window and
        verifies that a re-render completes correctly.
  \item \code{TestBasic}---iterates through every
        \code{RenderAlgorithm} entry, rendering the default view,
        View~5, and View~11 with both 32-bit and 64-bit iteration
        types.
  \item \code{TestReferenceSave}---saves and reloads reference orbits,
        verifying round-trip correctness.
  \item \code{TestVariedCompression}---tests reference-orbit
        compression at varying error thresholds.
  \item \code{TestImaginaLoad}---loads orbits stored in the Imagina
        format (\cref{subsec:imagina-format}).
  \item \code{TestStringConversion}---validates string-to-precision
        and precision-to-string conversion utilities.
  \item \code{TestPerturbedPerturb}---exercises perturbed iteration
        chains for numerical correctness.
  \item \code{TestGrowableVector}---validates the
        \code{GrowableVector} data structure
        (\cref{sec:disk_backed_growable_vectors}).
\end{itemize}
A separate \code{TestReallyHardView27} (not included in
\code{TestAll}) renders View~\#27, a period ${\sim}28$~billion point
requiring hours of computation; the test serves as an overnight
stress benchmark rather than a routine regression check.

% -----------------------------------------------------------------
\subsection{PNG Encoder and Save-Path Validation}
\label{subsec:png-validation}

The PNG encoder described in \cref{subsec:png-output} has a focused
hardware suite in \code{FractalSharkCudaTest}. The \code{CudaPng\_*}
cases encode identical input through the GPU encoder and the CPU's
current PNG save settings. LodePNG decodes
both outputs to RGBA16 for exact comparison with independently
serialized expected pixels. Coverage includes low channel bits,
transparency, padded strides, row and region boundaries, deterministic
output, noisy input, multiple IDAT chunks, workspace reuse, concurrent
streams, and invalid inputs.

Internal filtering tests reconstruct scanlines independently. Raw zlib
tests cover length/distance coding, overlapping copies, stored blocks,
checksums, and final-block termination. Huffman tests compare code
lengths and weighted cost with the bundled CPU implementation, including
degenerate alphabets, stable ties, and skewed frequencies requiring
length limiting. Header tests independently expand repetition codes.
Corruption tests separate PNG CRC failures from Adler-32 failures by
repairing the enclosing chunk CRC. Internal test interfaces in
\code{PngCompressionKernels.cuh} use production device routines and
also expose the fixed/stored baseline for size comparisons.

The \code{CudaPngSave\_*} cases exercise the production dispatch path:
both iteration widths, AA factors~1--4, odd dimensions, repeated saves
and resizing, copied/moved snapshots, render-pool frame publication,
synchronous Background completion, Unicode filenames, existing-file
skips, write failures, and switching back to queued CPU saves.
Comparisons with the CPU saver use settings where CPU and GPU coloring
agree, including zero palette rotation. The existing CPU save-pool and
PNG tests remain runnable without CUDA hardware.

Run the following PowerShell commands from the repository root after
building both Windows configurations:

\begin{lstlisting}[language={},upquote=true]
foreach ($configuration in 'Debug', 'Release') {
    & ".\$configuration\FractalSharkCudaTest.exe" --use-gpu `
        --filter=CudaPng_* --filter=CudaPngSave_*
    & ".\$configuration\FractalSharkTest.exe" `
        --filter=FractalSharkLib_SavePool* `
        --filter=FractalSharkLib_PngEncoding*
}
\end{lstlisting}

\paragraph{Existing render goldens.}
Two existing View0 cases exercise CPU64 and GPU1x32 PNG saves through
\code{SaveCurrentFractal}. Their fixtures drain the render pool, call
\code{CalcFractal(true)}, save the current frame, and wait for completion.
Only those cases are selected by the following commands:

\noindent\begin{minipage}{\linewidth}
\begin{lstlisting}[language={},upquote=true]
foreach ($configuration in 'Debug', 'Release') {
    & ".\$configuration\FractalSharkTest.exe" --use-gpu `
        --filter "RenderGolden_Basic_Cpu64_View0_Bits32_Store0_*" `
        --filter "RenderGolden_Basic_Gpu1x32_View0_Bits32_Store0_*"
}
\end{lstlisting}
\end{minipage}

Render goldens compare decoded RGBA16 pixels, so different compression
bytes alone require no golden update. Intentional pixel changes use the
selected-case update procedure in \cref{subsec:render-regression-goldens};
unaffected cases need not be regenerated or rerun.

\paragraph{Benchmark procedure.}
The standalone benchmark uses opaque uniform, gradient, seeded-noise,
and GPU-rendered Mandelbrot images at 3840~by~2160 resolution:

\begin{lstlisting}[language={},upquote=true]
.\Release\FractalSharkCudaTest.exe --use-gpu --filter=CudaPngBenchmark_4K
\end{lstlisting}

Run it after other builds and tests complete. Every output is decoded
and checked exactly. Compression size must not exceed the fixed/stored
baseline, and uniform, gradient, and Mandelbrot fixtures must improve
on that baseline. Timings are diagnostic rather than assertions;
\cref{subsec:png-benchmarks} defines the measurement boundaries and
reports the recorded results.

\paragraph{Compute Sanitizer on Windows.}
Select named-pipe transport before launching memcheck:

\begin{lstlisting}[language={},upquote=true]
$env:NV_COMPUTE_SANITIZER_LOCAL_CONNECTION_OVERRIDE = 'named-pipes'
& "$env:CUDA_PATH/compute-sanitizer/compute-sanitizer.exe" `
    --tool memcheck --error-exitcode 99 `
    .\Debug\FractalSharkCudaTest.exe --use-gpu `
    --filter=CudaPng_* --filter=CudaPngSave_*
\end{lstlisting}

The sanitizer communicates with its instrumented process. Named pipes
keep that Windows communication local without a TCP listener; TCP
transport can otherwise produce a firewall prompt attributed to the
test executable. Ordinary encoder tests start no network service.
The transport override is documented in NVIDIA's
\href{https://docs.nvidia.com/compute-sanitizer/ComputeSanitizer/index.html\#environment-variables}
{Compute Sanitizer environment-variable reference}.

% -----------------------------------------------------------------
\subsection{Selected Kernel Entry Points}
\label{sec:kernel-list}

This section lists representative CUDA entry points.  It is intentionally not
exhaustive: templated helpers, wrapper launch functions, generated
specializations, and internal kernels change more frequently than this appendix.

\begin{itemize}

\item \textbf{Mandelbrot base kernels} (\cref{sec:mandel-base-kernels})
  \begin{itemize}
    \item \code{mandel\_1x\_float}
    \item \code{mandel\_1x\_double}
    \item \code{mandel\_2x\_float}
    \item \code{mandel\_2x\_double}
    \item \code{mandel\_4x\_float}
    \item \code{mandel\_4x\_double}
    \item \code{mandel\_hdr\_float}
  \end{itemize}

\item \textbf{Mandelbrot perturbation kernels}
  \begin{itemize}
    \item \code{mandel\_1x\_double\_perturb\_bla}
    \item \code{mandel\_1xHDR\_float\_perturb\_bla} (\cref{sec:bilinear-approx})
    \item \code{mandel\_1xHDR\_float\_perturb\_lav2} (\cref{sec:la-v2-perturb,sec:perturb-only})
  \end{itemize}

\item \textbf{High-precision GPU reference kernels}
  (\cref{subsec:ref-orbit-gpu,subsec:fused-ntt-reference-step})
  \begin{itemize}
    \item \code{HpSharkReferenceSetupKernel}
    \item \code{HpSharkReferenceGpuLoop}
  \end{itemize}

\item \textbf{Utility kernels}
  \begin{itemize}
    \item \code{antialiasing\_kernel}
    \item \code{max\_reduce}
    \item \code{max\_kernel}
  \end{itemize}

\item \textbf{PNG encoding kernels} (\cref{subsec:png-output})
  \begin{itemize}
    \item Serialization and scanline filter selection
    \item Regional LZ77 encoding and dynamic Huffman construction
    \item Compaction, Adler-32, and PNG chunk assembly/CRC
  \end{itemize}

\item \textbf{Reference NTT primitives} (\cref{sec:ntt-multiply})
  \begin{itemize}
    \item \code{SharkNTT\_Primitives}
  \end{itemize}

\item \textbf{Disabled kernels}
  \begin{itemize}
    \item \code{mandel\_2x\_float\_perturb\_setup}
    \item \code{mandel\_2x\_float\_perturb}
  \end{itemize}

\end{itemize}

\subsubsection{Commentary}

\begin{itemize}

\item All the per-pixel CUDA kernels are in \code{gpu\_render.cu}. The two
most interesting are probably \code{mandel\_\allowbreak 1xHDR\_\allowbreak float\_\allowbreak perturb\_bla} and
\code{mandel\_\allowbreak 1xHDR\_\allowbreak float\_\allowbreak perturb\_lav2}, which are the ones I've spent the
most time on lately. For better performance at low zoom levels, consider
\code{mandel\_1x\_float\_perturb}, which leaves out linear approximation
and just does straight perturbation up to $\sim 10^{30}$, which corresponds
with the 32-bit float exponent range.

\item The \code{mandel\_1x\_float} (\cref{sec:mandel-1x-float}) is the classic 32-bit float Mandelbrot and 
is screaming fast on a GPU. This one is optimized with fused multiply-add for 
fun even though it is of limited practical use because one can barely zoom in before
pixellation sets in.

\end{itemize}

